\section{Introduction}

When practitioners choose a data attribution method to debug or understand their model, does it matter whether they're explaining BERT or GPT-2? Prior work has evaluated methods in isolation---EK-FAC on decoder-only LLMs \citep{grosse2023studying}, TRAK on BERT and CLIP separately \citep{park2023trak}---but no systematic comparison exists across architectures at matched conditions. We show that the answer depends critically on the method: TRAK is architecture-agnostic, TracIn favors encoders, and the expected EK-FAC advantage for decoders does not materialize.

Data attribution methods estimate the influence of training examples on model predictions. As foundation models grow in scale and deployment, understanding which training examples drive specific behaviors becomes essential for debugging, data quality assessment, and model auditing. Three prominent methods---TRAK, EK-FAC, and TracIn---make fundamentally different approximations to the intractable influence function, trading computational cost for accuracy. However, all prior evaluations have treated architecture as a fixed context rather than an experimental variable.

We hypothesize that attention structure should interact with attribution approximations. Encoder-only models like BERT use bidirectional attention, where each token attends to all others, creating dense gradient flow patterns. Decoder-only models like GPT-2 use causal attention masks that zero out half the attention weights, concentrating gradient flow in the lower triangle. These structural differences propagate through backpropagation, affecting the Hessian curvature that approximation methods must capture or circumvent.

To test this hypothesis, we conduct matched experiments comparing BERT-base and GPT-2 (both 12 layers, $\sim$110--125M parameters) on SST-2 mislabeled detection with 5\% injected label noise. We evaluate TRAK, EK-FAC, and TracIn at three compute budgets each, measuring mislabeled detection AUC---a standard influence function benchmark.

Our findings reveal three key insights:

First, \textbf{TRAK exhibits remarkable architecture-invariance}, with less than 1\% AUC difference between BERT and GPT-2 at all compute budgets. Its random projection mechanism appears to average out architecture-specific gradient patterns, providing robust cross-architecture attribution.

Second, \textbf{TracIn shows a directional BERT advantage} ($\sim$3\% at low compute). Its gradient dot-product approximation benefits from the dense bidirectional gradient flow in encoders.

Third, \textbf{EK-FAC does not show the predicted GPT-2 advantage}. Despite theoretical expectations that Kronecker factorization fits causal structures better \citep{grosse2023studying}, we observe only 0.27\% difference---statistically indistinguishable from zero. The Kronecker approximation appears more robust to attention structure than prior analysis suggested.

We validate the causal mechanism underlying these patterns. Attention sparsity differs by 98.8\% between architectures (BERT: 1.2\% upper-triangle zeros, GPT-2: 100\%). This structural difference creates an 11$\times$ gap in top Hessian eigenvalue (GPT-2: 0.502, BERT: 0.046), confirming that attention structure fundamentally shapes the loss landscape. All three attribution methods show $>$10\% relative AUC differences across architectures in at least one condition, demonstrating measurable architecture-method interaction.

These results provide actionable guidance for practitioners. For cross-architecture pipelines or when architecture may vary, TRAK is the robust default choice. For encoder-only work, TracIn may offer advantages. EK-FAC performs equivalently on both architectures, contrary to prior theoretical expectations. Our code and experiments are available for reproducibility.
